\chapter{Appendices}

\section*{Appendix A}\label{sec:appendixA}
In this part, we sketch a proof for Lemma \ref{lem:RK-local-truncation}.
\begin{proof}
We note that the second relation in equation \eqref{eqn:RK-fully-trunc} implies
\begin{equation*}
 u(t_{ni}) -   u(t_{n-1}) -\tau\sum_{j=1}^m a_{ij}  u_t^{nj} =
 \tau \sum_{j=1}^m a_{ij} (\dot u_*^{nj} - u_t(t_{nj}))+
 \eta_{ni}\quad \text{for} ~~i=1,2,\dots,m.
\end{equation*}
Then we substitute the first relation of \eqref{eqn:RK-fully-trunc} and derive that
for $i=1,2,\dots,m$
\begin{equation*}
 u(t_{ni}) -   u(t_{n-1}) -\tau\sum_{j=1}^m a_{ij}  u_t^{nj} =
 \tau \sum_{j=1}^m a_{ij}   \Big(\sum_{\ell=1}^k L_{\ell}(t_{n-1}+c_j\tau) f(u(t_{n-\ell})) - f(t_{nj})\Big)+
 \eta_{ni}.
\end{equation*}
Define $\tilde\eta_{ni}$ as the left hand side of the above relation. %equation \eqref{eqn:err_splitted}.
Now we apply Taylor's expansion at $t_{n-1}$  and use \eqref{eqn:stage-order} to derive
\begin{equation*}
\begin{aligned}
\tilde\eta_{ni} =& \sum_{l=1}^{m}\frac{\tau^l}{(l-1)!}\l(\frac{c_i^l}l -
\sum_{j=1}^ma_{ij}c_j^{l-1}\r)u^{(\ell)}(t_n) +
\frac1{m!}\int_{t_{n-1}}^{t_{ni}} (t_{ni}-s)^{m} u^{(m+1)}(s) \d s \\
&+ \frac\tau{(m-1)!}\sum_{j=1}^ma_{ij}\int_{t_{n-1}}^{t_{nj}}
(t_{nj}-s)^{m-1} u^{(m+1)}(s) \d s\\
&=\frac1{m!}\int_{t_{n-1}}^{t_{ni}} (t_n-s)^{m} u^{(m+1)}(s) \d s
+ \frac\tau{(m-1)!}\sum_{j=1}^ma_{ij}\int_{t_{n-1}}^{t_{nj}}
(t_{nj}-s)^{m-1} u^{(m+1)}(s) \d s
\end{aligned}
\end{equation*}
Then we obtain the estimate for $\tilde\eta_{ni}$, with $i=1,2,\dots,m$, that
\begin{equation*}
\|\tilde\eta_{ni}\|_{H^1(\Omega)} \leq C\tau^{m+1} \|u^{(m+1)}\|_{C([t_{n-1},t_n];H^1(\Omega))}.
\end{equation*}
This together with the approximation property of Lagrange interpolation lead to
\begin{equation*}
    \|\eta_{ni}\|_{H^1(\Omega)} \leq C\Big(\tau^{k+1} \| f(u)  \|_{C^{k}([t_{n-k},t_n]; H^1\II)}
    + \tau^{m+1} \|u\|_{C^{(m+1)}([t_{n-1},t_n];H^1(\Omega))} \Big). %, \quad \text{for} ~~i=1,2,\dots,m.
\end{equation*}
for $i=1,2,\dots,m$. Similarly, we have
\begin{equation*}
    u(t_{n}) -  u(t_{n-1}) -\tau\sum_{i=1}^m b_i  u_t^{ni} =
    \tau \sum_{i=1}^m b_i  \Big(\sum_{\ell=1}^k L_{\ell}(t_{n-1}+c_i\tau) f(u(t_{n-\ell})) - f(t_{ni})\Big)  +
    \eta_{n}.
\end{equation*}
Take the left hand side as $\tilde{\eta_n}$. Then Taylor expansion and \eqref{eqn:accuracy} imply
\begin{equation*}
\begin{aligned}
\tilde\eta_{n} =\frac1{p!}\int_{t_{n-1}}^{t_n} (t_n-s)^p u^{(p+1)}(s) \d s
+ \frac\tau{(p-1)!}\sum_{i=1}^mb_i\int_{t_{n-1}}^{t_{ni}}
(t_{ni}-s)^{p-1} u^{(p+1)}(s) \d s.
\end{aligned}
\end{equation*}
This together with the approximation property of Lagrange interpolation leads to
\begin{equation*}
     \|\eta_{ni}\|_{H^1(\Omega)} \leq C\Big(\tau^{k+1} \| f(u)  \|_{C^{k}([t_{n-k},t_n]; H^1\II)}
    + \tau^{p+1} \|u \|_{C^{p+1}([t_{n-1},t_n];H^1(\Omega))} \Big).
\end{equation*}
Using the choice that $k=\min(p, m+1)$, we derive the desired result.
\end{proof}



\section{List of IMEX-RK Methods}
In this part, we present some IMEX-RK schemes that satisfy
Assumptions (P1)-(P4).

\subsubsection{Example 1: first order IMEX}
This simple one-step case corresponds to the following IMEX scheme whose tableau reads:
\begin{table}[h]
    \centering
    \begin{tabular}[h]{c|c|c}
        $1$ & $1$ & $1$\\
        \hline
        ~ & $1$ & $1$
    \end{tabular}
    \caption{Butcher tableau for first order IMEX-RK}
\end{table}

This scheme agrees with Remark \ref{Remark:IMEX-Energy}. 
In this example,
$$\sigma(\lambda) = \frac{1}{1+\lambda}, \forall \lambda > 0.$$


\subsubsection{Example 2: second order IMEX}
The following Butcher tableau gives us a second order IMEXRK scheme:
\begin{table}[h]
    \centering
    \begin{tabular}[h]{c|cc|cc}
        $\gamma$ & $\gamma$ & ~ & $\gamma$ & ~\\
        $1$ & $1- \gamma$ & $\gamma$ & $\delta$ & $1-\delta$\\
        \hline
        ~ & $1- \gamma$ & $\gamma$ & $\delta$ & $1-\delta$
    \end{tabular}
    \caption{Butcher tableau for second order IMEX-RK}\label{tab:imex2}
\end{table}

in which $\gamma = 1+\frac{\sqrt{2}}2$, $\delta = 1-\frac1{2\gamma}$.
This scheme agrees with Remark \ref{Remark:IMEX-Energy}. 
In this example,
$$\sigma(\lambda) = \frac{1+(1+\sqrt2)\lambda}{(1+(1+\sqrt2/2)\lambda)^2}, \forall \lambda > 0$$

\subsubsection{Example 3: third order IMEX}
The following Butcher tableau gives us a third order IMEXRK scheme:
\begin{table}[h]
    \centering
    \begin{tabular}[h]{c|cccc}
        $3/5$ & $3/5$  \\
        $3/2$ & $15/32$ & $33/32$    \\
        $19/20$ & $2/5$ & $-357/640$  &$709/640$ \\
        $1$ & $2825/756$ & $-232/297$  &$-6400/231$ & $103/4$ \\
        \hline
        ~ &  $2825/756$ & $-232/297$  &$-6400/231$ & $103/4$ \\
    \end{tabular}
    \begin{tabular}[h]{c|cccc}
        $3/5$ & $3/5$ \\
        $3/2$ &  $51/64$ & $45/64$  \\
        $19/20$ &  $2/5$ & $4841/11520$  & $299/2304$ \\
        $1$ &  $103/342$ & $125/378$  & $-26/297$ & $2000/4389$\\
        \hline
        ~ &  $103/342$ & $125/378$  & $-26/297$ & $2000/4389$\\
    \end{tabular}
    \caption{Butcher tableau for third order IMEX-RK}
\end{table}

This scheme agrees with Remark \ref{Remark:IMEX-Energy}. 
In this example,
$$\sigma(\lambda) = \frac{1228800+ 33778560\lambda + 55256268 \lambda^2 + 5250325\lambda^3 }
{1228800+ 35007360\lambda + 89649228 \lambda^2 + 77600673\lambda^3 +21689019\lambda^4}, \forall \lambda > 0$$

\subsubsection{Example 4: forth order IMEX scheme for linear problem}

 The following Butcher tableau gives us an IMEXRK scheme 
 which is fourth order for linear problem and third order for semilinear problem:
 \linespread{1.5}
 \begin{table}[h]
     \centering
     \begin{tabular}[h]{c|cccccc}
         $\frac{1}{5}$ & $\frac{1}{5}$ \\
         $\frac{2}{5}$ & $\frac{1}{5}$&$\frac{1}{5}$\\
         $\frac{3}{5}$ & $\frac{7}{5}$&$-1$&$\frac{1}{5}$\\
         $\frac{4}{5}$ & $\frac{3989}{1025}$&$-\frac{3343}{1025}$&$-\frac{31}{1025}$&$\frac{1}{5}$\\
         $1$ & $\frac{79}{120}$&$-\frac{91}{120}$&$\frac{149}{120}$&$-\frac{41}{120}$&$\frac{1}{5}$\\
         $1$ & $\frac{79}{120}$&$-\frac{91}{120}$&$\frac{149}{120}$&$-\frac{41}{120}$&$0$ &$\frac{1}{5}$\\
         \hline
         ~ & $\frac{79}{120}$&$-\frac{91}{120}$&$\frac{149}{120}$&$-\frac{41}{120}$&$0$ &$\frac{1}{5}$\\
     \end{tabular}
     \begin{tabular}[h]{c|cccccc}
         $\frac{1}{5}$ & $\frac{1}{5}$ \\
         $\frac{2}{5}$ & $\frac{5299}{30120}$&$\frac{1357}{6042}$\\
         $\frac{3}{5}$ & $\frac{8681383}{58516770}$&$\frac{5413353}{9752795}$&$-\frac{6051439}{58516770}$\\
         $\frac{4}{5}$ & $0$&$\frac{25905769}{16101930}$&$-\frac{9808291}{16101930}$&$-\frac{10113}{50635}$\\
         $1$&$0$ & $\frac{79}{120}$&$-\frac{55873}{32520}$&$\frac{96299}{32520}$&$-\frac{5863}{6504}$\\
         $1$ &$0$ &$\frac{79}{120}$&$-\frac{91}{120}$&$\frac{149}{120}$&$-\frac{41}{120}$ &$\frac{1}{5}$\\
         \hline
         ~ &$0$ &$\frac{79}{120}$&$-\frac{91}{120}$&$\frac{149}{120}$&$-\frac{41}{120}$ &$\frac{1}{5}$\\
     \end{tabular}
     \caption{{Butcher tableau for fourth order IMEX-RK}}
 \end{table}

% % 0.199230027598864
% % 0.332262474162124
% % 0.461570781849613
% % 0.157095284449734
% % -0.150158568060334
This scheme does not agree with Remark \ref{Remark:IMEX-Energy}. 
 In this example,
 $$\sigma(\lambda) = \frac{75000-7500 \lambda^2 +1000\lambda^3 +225 \lambda^4}
 {24(5+\lambda)^5}, \forall \lambda > 0.$$
Although there are some negative coefficients in the numerator, it is still positive for all $\lambda \geq 0$.

\subsubsection{Example 5: forth order IMEX scheme for linear and semilinear problem}

 The following Butcher tableau gives us an IMEXRK scheme 
 which is fourth order for linear and semilinear problem:
 \linespread{1.5}
 \begin{table}[h] \label{Tab:IMEX47}
     \centering
     \begin{tabular}[h]{c|ccccccc}
         $\frac{1}{7}$ & $\frac{1}{7}$ \\
         $\frac{2}{7}$ & $\frac{1}{7}$&$\frac{1}{7}$\\
         $\frac{11}{25}$ & $\frac{1354}{4375}$&$-\frac{104}{4375}$&$\frac{27}{175}$\\
         $\frac{4}{5}$ & $a_{41}$&$a_{42}$&$a_{43}$&$\frac{9}{25}$\\
         $\frac{13}{15}$ & $a_{51}$&$a_{52}$&$a_{53}$&$1$& $\frac{1}{15}$\\
         $\frac{14}{15}$ & $-\frac{14605}{25536}$&$1$&$\frac{1}{2}$&$\frac{1}{3}$&$-\frac{50399}{127680}$&$\frac{1}{15}$\\
         $1$&$a_{71}$&$a_{72}$&$a_{73}$&$a_{74}$&$0$& $\frac{1}{15}$& $\frac{1}{15}$\\
\hline
        ~&$a_{71}$&$a_{72}$&$a_{73}$&$a_{74}$&$0$& $\frac{1}{15}$& $\frac{1}{15}$\\
     \end{tabular}
     \begin{tabular}[h]{c|ccccccc}
        $\frac{1}{7}$ & $\frac{1}{7}$ \\
        $\frac{2}{7}$ & $\frac{1}{7}$&$\frac{1}{7}$\\
        $\frac{11}{25}$ & $-\frac{3229}{8750}$&$1$&$-\frac{1671}{8750}$\\
        $\frac{4}{5}$  & $-1$ & $\hata_{41}$&$\hata_{42}$&$\hata_{43}$\\
        $\frac{13}{15}$ & $1$ & $\hata_{51}$&$\hata_{52}$&$\hata_{53}$&$\hata_{54}$\\
        $\frac{14}{15}$ & $1$ & $\frac12$&$\frac13$&$\hata_{63}$&$\hata_{64}$&$\hata_{65}$\\
        $1$&$\hata_{70}$&$\hata_{71}$&$\hata_{72}$&$\hata_{73}$&$\hata_{74}$&$\hata_{75}$& $\hata_{76}$\\
\hline
~&$\hata_{70}$&$\hata_{71}$&$\hata_{72}$&$\hata_{73}$&$\hata_{74}$&$\hata_{75}$& $\hata_{76}$\\
\end{tabular}
     \caption{{Butcher tableau for fourth order IMEX-RK}}
 \end{table}
 \begin{table}[h] \label{Tab:coef}
    \centering
    \begin{tabular}[h]{|c|c|}
\hline
$a_{41}$ & $24058589213/79084733000$ \\
$a_{42}$ & $299597908/760430125$ \\
$a_{43}$ & $-23336559/90382552$ \\ \hline
$a_{51}$ & $-419174425120649204186119/79218920548212919110000$ \\
$a_{52}$ & $477902525326525491346619/41132901053879784922500$ \\
$a_{53}$ & $-558440160387285872007719/85556434192069952638800$ \\ \hline
$a_{71}$ & $4922197/9687600$ \\
$a_{72}$ & $-2268847/3936600$ \\
$a_{73}$ & $3238175/4094064$ \\
$a_{73}$ & $869063/6036120$ \\\hline
$\hata_{41}$ & $55740853326839983334621/27520768939834456977600$ \\
$\hata_{42}$ & $-11060173572755944147541/14289630026452506507600$ \\
$\hata_{43}$ & $2329344893911909455707/4246061493574459076544$ \\ \hline
$\hata_{51}$ & \scriptsize $-71186087675929044223656193931944384573278473/3438510885083184524462404153200164172975000$ \\
$\hata_{52}$ & \scriptsize $14268306999047363197901818351205461479990553/399215549486485350905095178799636834225000$ \\
$\hata_{53}$ & \scriptsize $-164196843957237023227423093901438282094053/11625156801046453418356371606645424612632$ \\
$\hata_{54}$ &  $-2777877807409624755035963389/2652337211701823798398155000$ \\ \hline
$\hata_{63}$ &  $-1478910748534537278959/506708836650241118208$ \\
$\hata_{64}$ &  $29121578321788444001/11310465103800024960$ \\
$\hata_{65}$ &  $-10436070861733408099/18766993950008930304$ \\ \hline
$\hata_{70}$ &  $-12432250836521654063395/85826994234468729923442$ \\
$\hata_{71}$ &  $8194496564065915317284063497/9238417659398214088959296880$ \\
$\hata_{72}$ &  $-2996905844557888599918137227/3754072727815662246851353080$ \\
$\hata_{73}$ &  $13920174158309482913767468375/19521178184641443683627036016$ \\
$\hata_{74}$ &  $1389846140555607156290661143/5756244849317348778505408056$ \\ 
$\hata_{75}$ &  $84635598876295853750/14304499039078121653907$ \\ 
$\hata_{76}$ &  $1369093183594155319761/14304499039078121653907$ \\ \hline
\end{tabular}
    \caption{related coefficients in Table \ref{Tab:IMEX47}}
\end{table}

Related coefficients are listed in Table \ref{Tab:coef}. 
This scheme does not agree with Remark \ref{Remark:IMEX-Energy}. 
 In this example,
\begin{align*}
\sigma(\lambda) = &-\Big(- 18533185137819\lambda^5 + 245682733504208\lambda^4 - 1917903570331840\lambda^3 \\
&+ 4891758175886400\lambda^2 + 9437315024592000\lambda - 61167782566800000\Big)\\
&/\Big(14794928512(9\lambda + 25)(\lambda + 7)^2(\lambda + 15)^3\Big) , \forall \lambda > 0.
\end{align*}
Although there are some negative coefficients in the numerator, it is still positive for all $\lambda \geq 0$.

\subsubsection{Find IMEX--RK Scheme}

In this part we make a short introduction to how we find the table of IMEX-RK.

Searching the table is based on undetermined coefficient method.
To this end, our purpose is to find a $k$-th order scheme.
First we need to know is that, for a $m$-stage IMEX--RK table, the total degree of freedoms is about $m^2$,
but the requirements of stage order is about $2^k$,
because the replacement of $b, \hatb$ and $A, \hatA$ should keep the requirements,
which means that the number of stages will be far more than the scheme order.

However, when we focus on the implicit table $A$ itself, the scheme is reduced to a diagonal implicit Runge--Kutta (DIRK)
scheme, and the replacement conditions of $b, \hatb$ and $A, \hatA$ is no need to be kept, so the total number of requirements is now algebraic to $k$.
We can fix a small $m$ and find a small $A$ to the DIRK scheme.This is not hard.

In fact, we have plenty of degree of freedoms during finding the DIRK table,
so we can fix $c_i$ first because it is high order in the stage order requirements.
We can also make the diagonal terms in $A$ same, so that the solvers are same in each stages.
If we have other requirements to $A$, like $\sigma(\lambda)>0$,
we can search among the rest degree of freedoms.

After a first $A$ is generated in a small size, we will solve $\hatA$ now.
The current number of stages may be not enough for meeting all the stages orders,
so we should add some stages.
We do not need to hesitate whether too much stages are added, because that the number of stages is always far more than the scheme order for high order schemes.

The adding of stages is just repeat the last stage in the implicit table $A$, like what we have done in the above forth order scheme.
We can see that the last two stages of the implicit table is almost the same, with only an extra $0$.
With focus on DIRK scheme, this has nothing change and the stage order is also kept,
but it gives us extra degree of freedoms so we can solve a larger system.
One can also find that the $b$ and $\hatb$ are also similar.
In fact $b_\sigma$ and $\hatb_\sigma$ are also same except the swap of last two terms.
Either of the last two terms refers to $c_i=1$ so the swapping works.

Here the last and the hardest task is search $\hata$, this can be accessed by undetermined coefficient method.
We should add an extra stage if the equation is not solvable, the degree of freedom is not enough,
and an extra stage, again and again until we can find the table.
We guess perhaps there exist arbitrary high order IMEX-RK schemes.

So the search has the following steps:
\begin{itemize}
    \item[{\bf Step 1:}] Give $c$ manually.
    \item[{\bf Step 2:}] Solve $b$ and $A$ for DIRK scheme with a small number of stages $m$. For requirements to $A$, do it here.
    \item[{\bf Step 3:}] Calculate $\hatb$ or keep $b$ as $\hatb$.
    \item[{\bf Step 4:}] Try to solve $\hatA$ from the stage order requirements with undetermined coefficient method. Output the tables.
    \item[{\bf Step 5:}] If the equations are not solvable, add $m$, go to Step 3.
\end{itemize}




\section{Proof of Parareal Stability to Runge--Kutta Methods}
In this part, we provide detailed proof of proposition \ref{prop:lobatto-34}.

\subsection*{Proof for the 3-stage Lobatto IIIC scheme \eqref{3-stage}}
Similar to the proof of proposition \ref{prop:lobatto-2}, we aim to show that for any $J\ge 2$
\begin{equation}\label{eqn:lobatto-3}
    \sup_{s\in(0,\infty)}  \left|\frac{(1+s) r(-s/J)^J-1}{s}\right| \le 0.31,
    \quad
    \text{where}~~ r(-s) = \frac{24-6s}{s^3 + 6s^2 + 18s + 24}.
\end{equation}
Letting $\alpha = 0.69$ and $\beta=1.02$,
Lemmas \ref{lem:exp} and \ref{lem:exp-102} implies that
$\kappa_\alpha \le 1-0.69 = 0.31$ and $\kappa_\beta \le 0.31$, respectively.

Next, we show the claim that $e^{-\beta s} \le r(-s) \le e^{-\alpha s}$ for all $s\in(0,s_*)$ with $s_*=2$. To begin with, we shall prove that $r(-s) > e^{-\beta s}$ for all $s\in(0,\infty)$, which is equivalent to
$$\psi(s) = e^{\beta s} (-6s+24) - (s^3+6s^2+18s+24) >0
\quad \forall ~s\in (0, s_*).$$
We note that
$$\psi^{(4)}(s) = 6\beta^3e^{\beta s}(-\beta s + 4\beta - 4).$$
has a unique root in $(0,\infty)$, namely
$s_0= \frac{4\beta - 4}{\beta} \approx0.0784$, and hence
$\psi^{(4)}(s)> 0$ for all $s\in(0,s_0)$.
Besides, we observe that
$\psi^{(3)}(0) = 0.742 >0$, $\psi^{(3)}(s_0) \approx 0.762 >0$.
Therefore $\psi^{(3)}(s)$ has a unique root
in $(s_0, +\infty)$,
denoted as $s_1$. By means of the fixed point iteration, we know that
$s_1\approx 0.4832$.
Similarly, since $\psi''(0) =  0.730>0$ and $\psi''(s_1) =  1.00>0$, we conclude that $\psi''(s)$ is always positive in $[0,s_1]$ and it
has a unique root $s_2\in (s_1,\infty)$, and we find $s_2\approx 0.9980$.
Repeating the argument, we are able to show that
$\psi'(s)$ keeps positive in $[0,s_2]$ and
the unique root in $(s_2,\infty)$ locates at
$s_3 \approx 1.5344$.
Finally, we observe that $\psi(0) = 0$ and $\psi(s_3) \approx 1.401>0$, so $\psi$ is positive in $(0,s_3]$
and it admits a unique root at $s_4\in(s_3,\infty)$.
Noting that $\psi(s_*) \approx 0.2873 > 0$,
we conclude that $r(-s) > e^{-\beta s}$ for all $s\in(0,s_0)$.

Next we will show that $r(-s) < e^{-\alpha s}$ in $(0, s_*)$, which is equivalent to show
\begin{equation*}
    \varphi(s) = e^{\alpha s} (-6s+24) - (s^3+6s^2+18s+24) < 0 \quad \text{for all} ~s\in (0, s_*).
\end{equation*}
We note the fact that
\begin{equation*}
\varphi^{(4)}(s) = 6\alpha^3e^{\alpha s}(-\alpha s + 4\alpha - 4)<0\quad \text{for all} ~s\in (0, \infty).
\end{equation*}
Meanwhile, we have $\varphi^{(4)}(0)<0$, $\varphi^{(3)}(0)<0$, $\varphi''(0)<0$, $\varphi'(0)<0$, and $\varphi(0)<0$.
Those together imply $\varphi(s) <0$ for any
$s\in (0, \infty)$. Therefore $r(-s) < e^{-\alpha s}$.

As a result, for any $J\ge 2$, we arrive at $e^{-\beta s} \le r(-s/J)^J \le e^{-\alpha s}$
for all $s\in(0, J s_*)$, that further implies the estimate
\begin{equation*}
    \sup_{s\in(0, J s_*)}
    \left|\frac{(1+s) r(-s/J)^J-1}{s}\right| \le 0.31.
\end{equation*}

Next, we aim to prove the claim that
\begin{equation}\label{claim:bound-3}
   \sup\limits_{(s_*, \infty)}
   \big|r(-s)\big| = \sup\limits_{(s_*, \infty)}
   \Big|\frac{24-6s}{s^3 + 6s^2 + 18s + 24}\Big|
   \le  0.15.
\end{equation}
To begin with, we show that $\frac{\d}{\d s}r(-s)$ admits a unique root in $(2, \infty)$. % ,  denoted as $s_5$.
We note that
$$\frac{\d}{\d s}r(-s) = \frac{12(s^3-3s^2-24s-48)}
{(s^3 + 6s^2 + 18s + 24)^2}.$$
and hence it is sufficient to show that
$\eta(s) = s^3-3s^2-24s-48$
has a unique root in $(2, \infty)$.
Since $\eta'(s)$ has two roots, $-2$ and $4$,
and $\eta(-2)=-20<0$, $\eta(4)=-128<0$,
we conclude that $\eta(s) <0$ for all $s\in[-2,4]$, and
$\eta(s)$ admits a unique root in $(4, \infty)$,
namely $s_5 \approx 7.235$. %7.234531027367
Therefore $r(-s)$ is decreasing in $(s_*,s_5)$
and increasing in $[s_5,\infty)$.
Noting that fact that $r(-s_*) \approx 0.130$,
$r(-s_5)\approx -0.0229$ and $r(-\infty)=0$, we obtian
$
    \sup\limits_{(s_*, \infty)}
    \big|r(-s)\big| = \big|r(-s_*)\big|
    \le  0.15.
$

Therefore, we derive for $s\in (Js_*,\infty)$ and $J\ge 2$
\begin{equation*}
\begin{split}
    \left|\frac{(1+s) r(-s/J)^J-1}{s}\right|
    &\le \frac{1+s}{s}|r(-s/J)^J|+s^{-1}
    \le \frac{1+Js_*}{Js_*} (0.15)^J + (Js_*)^{-1}    \\
    &\le  \frac{5}{4} (0.02)^2 + 4^{-1}
    \approx 0.251 \le 0.31.
\end{split}
\end{equation*}
This completes the proof of \eqref{eqn:lobatto-3} as well as the proposition.
\qed

\subsection*{Proof for the 4-stage Lobatto IIIC scheme \eqref{4-stage}}
Similar to the proof of Proposition \ref{prop:lobatto-2}, we aim to show that for any $J\ge 2$
\begin{equation}\label{eqn:lobatto-4}
    \sup_{s\in(0,\infty)}  \left|\frac{(1+s) r(-s/J)^J-1}{s}\right| \le 0.31,
    ~~
    \text{where}~~ r(-s) =
    \frac{12s^2-120s+360}{s^4+12s^3+72s^2 + 240s + 360}.
\end{equation}
Letting $\alpha = 0.69$ and $\beta=1.02$,
Lemmas \ref{lem:exp} and \ref{lem:exp-102} implies that
$\kappa_\alpha \le 1-0.69 = 0.31$ and $\kappa_\beta \le 0.31$, respectively. Next, we show the claim that
\begin{equation}\label{claim:sand}
    e^{-\beta s} \le r(-s) \le e^{-\alpha s}~~ \text{for all}~~s\in(0,s_*)~~ \text{with}~s_*=6.8.
\end{equation}

To begin with, we show that, for $s\in(0,\infty)$, $r(-s) \ge e^{-\beta s}$, which is equivalent to
$$\psi(s) = (12s^2-120s+360)e^{\beta s} -(s^4+12s^3+72s^2 + 240s + 360)
 \ge 0.$$
Define $h(s) = \beta^2s^2 + (10\beta-10\beta^2)s+(30\beta^2-50\beta+20)$,
then we have
$$\psi^{(5)}(s) = 12\beta^3e^{\beta s}
\left[\beta^2s^2 + (10\beta-10\beta^2)s+(30\beta^2-50\beta+20)\right]
=12\beta^3e^{\beta s}h(s).$$
Here $h(s)$ is a quadratic polynomial, whose
minimum locates at $\frac{5\beta-5}{\beta}$.
Therefore $h(s) \geq h (\tfrac{5\beta-5}{\beta} ) > 0$.
Then $\psi^{(5)}(s) = 12\beta^2e^{\beta x}h(s)-24 > 12\times 2-24>0$.
Meanwhile, simple computation yields
$$ \psi(0)=0 \qquad \text{and}
\qquad \psi^{(k)}(0)>0~~\quad \text{with}~~ \quad
1\le k\le 5.  $$
Then we conclude that  $\psi(s) > 0$ for all  $s\in (0, \infty)$, and hence
 $r(-s) \ge e^{-\beta s}$ in $(0, \infty)$.

Next we show the bound that $r(s) \le e^{-\alpha s}$ for $s\in (0,s_*)$, which is equivalent to show
$$\varphi(s) = (12s^2-120s+360)e^{\alpha s}
-(s^4+12s^3+72s^2 + 240s + 360) \le 0 \qquad \forall~s\in (0,s_*).$$
Similar to the preceding argument, let $g(s) = (20-50 \alpha + 30 \alpha^2  +10\alpha s -10\alpha^2s +\alpha^2s^2)$. Then
$$  \varphi^{(5)}(s) = 12\alpha^3e^{\alpha s}
(20-50 \alpha + 30 \alpha^2
+10\alpha s -10\alpha^2s +\alpha^2s^2)
= 12\alpha^3e^{\alpha s}  g(s). $$
Here $g$ is a quadratic polynomial with minimum at $\frac{5\alpha-5}{\alpha}$. Therefore, $g(s)\ge g(\tfrac{5\alpha-5}{\alpha}) \approx -2.62$.
Meanwhile, we observe that $g(0) = -0.217 <0$, so
there is a unique root of $g$ in $(0, \infty)$.
It is easy to find that, by means of the fixed point iteration,
that root locates at $s_0\approx0.0993$. Then $\varphi^{(5)}(s) \le 0$
for all $s\in [0,s_0]$ and $\varphi^{(5)}(s) \ge 0$ for $s\in (s_0,\infty)$. Noting that $\varphi^{(4)}(s_0) < 0$ and $\varphi^{(4)}(0) < 0$, so $\varphi^{(4)}(s) < 0$ in $[0,s_0]$ and $\varphi^{(4)}$ admits a unique root in $(s_0,\infty)$, named as $s_1$.
Then the fixed point iteration implies $s_1\approx  1.6849$.
Repeating this argument, we are able to show that
$\varphi^{(3)}(s) < 0$ in $[0,s_1]$ and $\varphi^{(3)}$ has a unique root $(s_1,\infty)$, namely $s_2\approx 3.0558$.
Then we derive that
$\varphi{''}(s) < 0$ in $[0,s_2]$ and  $\varphi{''}(s)$
has a unique root in $(s_2,\infty)$,
denoted as $s_3\approx 4.3640$.
Similarly, $\varphi{'}(s) < 0$ in $[0,s_3]$ and  $\varphi{'}(s)$
has a unique root in $(s_3,\infty)$, named as $s_4\approx 5.6285$.
Finally, since $\psi(0)=0$ and $\varphi(s_4)<0$, we conclude that
$\varphi(s)<0$ in $(0,s_4]$, and $\varphi(s)$ has a unique root in $(s_4,\infty)$. Then the fact that $\varphi(s_*) \approx -447.65<0$ implies $\varphi(s)<0$ in $(0,s_*)$.
This completes the proof of the claim \eqref{claim:sand}.
As a result, for any $J\ge 2$, we arrive at $e^{-\beta s} \le r(-s/J)^J \le e^{-\alpha s}$
for all $s\in(0, J s_*)$ which  implies
\begin{equation*}
    \sup_{s\in(0, J s_*)}
    \left|\frac{(1+s) r(-s/J)^J-1}{s}\right| \le 0.31.
\end{equation*}

Next, we intend to show the claim that
\begin{equation}\label{claim:bound-4}
 \sup\limits_{(s_*, \infty)}
   \big|r(-s)\big| =   \sup\limits_{(s_*, \infty)}  \big|\frac{12s^2-120s+360}{s^4+12s^3+72s^2 + 240s + 360}\Big| \le  0.02.
\end{equation}
In order to establish a bound for the supremum, we note
$$\frac{d}{ds}r(-s)
= \frac{-24s^5+216s^4+1440s^3-1440s^2-43200s-129600}
{(s^4+12s^3+72s^2 + 240s + 360)^2}, $$
and we will show that it admits a unique root in $(s_*, \infty)$,
denoted by $s_5$. Noting that, with
$\mu(s) = -s^5+9s^4+60s^3-60s^2-1800s-5400, $
we have
$$\frac{d}{ds}r(-s) = \frac{24\mu(s)}{(s^4+12s^3+72s^2 + 240s + 360)^2}.$$
So it suffices to show that $\mu(s)$ admits a unique  root in $(s_*, \infty)$. Since
$\mu^{(3)}(s)=-60s^2+216s+360,$ being a quadratic polynomial,
it gains the maximum at $1.8$ where
$\mu^{(3)}(1.8) = 554.4>0$.
Noting that $s_*>1.8$ and $\mu^{(3)}(s_*)=-945.6<0$,
we conclude that $\mu^{(3)}(s)<0$ for all $s\in(s_*, \infty)$.
Moreover, since $\mu''(s_*) \approx 1033.3 >0$ and $\mu''(\infty) =-\infty$, $\mu''(s)$ admits a unique root $s_6\approx 7.6503 \in (s_*, \infty)$. Then we know that $\mu''(s) > 0$ in $(s_*, s_6)$
and $\mu''(s) < 0$ in $(s_6, \infty)$.
Similarly, we have $\mu'(s_*) >0$ and $\mu'(s_6) >0$,
so  $\mu'(s) > 0$ in $[s_*, s_6]$. This together with the fact  $\mu'(\infty) <0$ implies that
$\mu'(s)$ has a unique root $s_7\approx 10.166 \in (s_6, \infty)$.
Finally, using the facts that $\mu(s_*)>0$ and $\mu(s_7)>0$,
we know $\mu(s) >0 $ in  $(s_*, s_7)$ and
 $\mu(s)$ has a unique root $s_8 \approx 12.28 \in (s_7, \infty)$.
Therefore $r(-s)$ is increasing in $(s_*, s_8)$,
and decreasing in $(s_8, \infty)$.
Noting that $r(-s)\approx 0.0088$, $r(-s_8)\approx 0.0118$, and $r(-\infty)=0$, we arrive at
$\sup\limits_{(s_*, \infty)}
\big|r(-s)\big| \le 0.02.$


As a result, the estimate \eqref{claim:bound-4} implies that
for $s\in (Js_*,\infty)$ and $J\ge 2$
\begin{equation*}
\begin{split}
    \left|\frac{(1+s) r(-s/J)^J-1}{s}\right|
    &\le \frac{1+s}{s}|r(-s/J)^J|+s^{-1}
    \le \frac{1+Js_*}{Js_*} (0.02)^J + (Js_*)^{-1}    \\
    &\le  \frac{14.6}{13.6} (0.02)^2 + 13.6^{-1}
    \approx 0.074 \le 0.31.
\end{split}
\end{equation*}
This completes the proof of \eqref{eqn:lobatto-4} as well as the proposition.
\qed



% \bibliographystyle{siam}

